\chapter*{Summary}
\setheader{Preface}

This literature study presents an extensive overview of lifelong learning techniques.
Origami robots, characterized by their construction using flexible materials and origami techniques, exhibit notable compliance and adaptability.
However, these traits also lead to a faster rate of degradation than traditional robots, necessitating the development of adaptive lifelong learning strategies for their effective and sustained operation.

Central to the study is the exploration of the challenges inherent in lifelong learning for robotics, notably catastrophic forgetting, efficient memory handling, and the detection of shifts in the robot's dynamics.
These challenges underscore the importance of a robot's ability to retain previously learned knowledge, efficiently store and retrieve memories, and recognize significant changes in its operational environment or in its own mechanical dynamics.

The study delves deeply into various machine-learning-based methods for managing the evolving dynamics of robots due to wear and tear.
It examines `blackbox' methods, which, despite their lack of interpretability, offer significant benefits in terms of performance for lifelong learning tasks.
These methods are scrutinized for their efficacy in developing control policies and dynamic models that enable robots to adapt continually to changing conditions.

Further, the study explores dual architectures in lifelong learning, which consist of separate modules tackling different aspects of the challenge.
This includes policy/control modules combined with knowledge bases, as well as multiple control policies designed to handle varying levels of uncertainty.
The composition and maintenance of knowledge bases are also addressed, emphasizing the importance of the type and form of knowledge stored, its retrieval, and the strategies for keeping the knowledge base relevant and up-to-date throughout the robot's lifespan.

In summary, this literature study provides a foundational perspective on the application of lifelong learning in the field of origami robotics.
It highlights the necessity and challenges of adaptive learning in this context and outlines potential machine learning and control strategies to address these challenges, thus laying the groundwork for further research and development in adaptive control and learning for soft, flexible robots.

\begin{flushright}
{\makeatletter\itshape
    \@author \\
    Delft, November 2023
\makeatother}
\end{flushright}
